\documentclass[a4,11pt]{aleph-notas}
% Se puede ver la documentación aquí:
% https://github.com/alephsub0/LaTeX_aleph-notas

% -- Paquetes adicionales
\usepackage{enumitem}
\usepackage{url}
\usepackage{array}
\usepackage{booktabs}
\usepackage{longtable}
\hypersetup{
    urlcolor=blue,
    linkcolor=blue,
}

% -- Datos
\institucion{Facultad de Ciencias Exactas, Naturales y Ambientales}
\carrera{Catálogo STEM}
\asignatura{Álgebra lineal}
\tema{Clase 07: Independencia lineal, bases y dimensión}
\autor{Andrés Merino}
\fecha{Periodo 2026-01}

\logouno[0.16\textwidth]{Logos/logoPUCE_04_ac}
\definecolor{colortext}{HTML}{1957d1}
\definecolor{colordef}{HTML}{1957d1}
\fuente{montserrat}

\begin{document}

\encabezado 

%%%%%%%%%%%%%%%%%%%%%%%%%%%%%%%%%%%%%%%%
\section*{Resultado de Aprendizaje}
%%%%%%%%%%%%%%%%%%%%%%%%%%%%%%%%%%%%%%%%

%%%%%%%%%%%%%%%%%%%%%%%%%%%%%%%%%%%%%%%%
\subsection*{RdA de la asignatura:}
\begin{itemize}[leftmargin=*]
    \item \textbf{RdA 1:} Comprender los conceptos fundamentales del Álgebra Lineal, incluyendo el estudio de matrices, determinantes, sistemas de ecuaciones lineales y espacios vectoriales, destacando su importancia en el análisis y resolución de problemas matemáticos.
    \item \textbf{RdA 2:} Resolver problemas prácticos y teóricos relacionados con el Álgebra Lineal, utilizando técnicas de cálculo con matrices, determinantes y sistemas de ecuaciones, así como en el análisis de espacios vectoriales.    
    % \item \textbf{RdA 3:} Aplicar los principios y métodos del Álgebra Lineal en una variedad de contextos demostrando comprensión de su relevancia y utilidad en situaciones prácticas y teóricas.
\end{itemize}


%%%%%%%%%%%%%%%%%%%%%%%%%%%%%%%%%%%%%%%%
\subsection*{RdA de la actividad:}
\begin{itemize}[leftmargin=*]
    \item Reconocer cuándo un conjunto de vectores es linealmente dependiente o independiente.
    \item Determinar si un conjunto de vectores genera un espacio vectorial.
    \item Interpretar la noción de base y dimensión de un espacio vectorial.
\end{itemize}

%%%%%%%%%%%%%%%%%%%%%%%%%%%%%%%%%%%%%%%%
\section*{Introducción}
%%%%%%%%%%%%%%%%%%%%%%%%%%%%%%%%%%%%%%%%

%%%%%%%%%%%%%%%%%%%%%%%%%%%%%%%%%%%%%%%%
\paragraph{Control de avance:}
Antes de iniciar la clase, se aplicará el control de avance con un ejercicio breve sobre combinación lineal y otro sobre identificación de dependencia lineal.

\paragraph{Pregunta inicial:}
¿Cómo distinguimos entre un conjunto que genera un espacio y una base de ese mismo espacio?

%%%%%%%%%%%%%%%%%%%%%%%%%%%%%%%%%%%%%%%%
\section*{Desarrollo}
%%%%%%%%%%%%%%%%%%%%%%%%%%%%%%%%%%%%%%%%

%%%%%%%%%%%%%%%%%%%%%%%%%%%%%%%%%%%%%%%%
\subsection*{Actividad 1: Espacio generado e independencia lineal}

La clase inicia con una recuperación breve de combinaciones lineales y del papel del campo de escalares. Luego, mediante clase magistral y práctica guiada, se introduce el espacio generado por un conjunto de vectores, se estudian criterios de independencia lineal y se discute cómo reconocer cuándo un conjunto genera un espacio vectorial.

\paragraph{¿Cómo lo haremos?}
\begin{itemize}[leftmargin=*]
    \item \textbf{Clase magistral:} se presentan las definiciones de conjunto generador, independencia lineal, base y dimensión. Se usan ejemplos en $\mathbb{R}^n$, matrices, funciones y polinomios. Se usa el resumen \href{https://fcena-puce.github.io/AlgebraLineal-05-N0048/01Resumen/Resumen07.pdf}{Resumen07.pdf}.
    \item \textbf{Resolución de ejercicios:} se guiará a los estudiantes en la solución de ejercicios del documento \href{https://fcena-puce.github.io/AlgebraLineal-05-N0048/04Ejercicios/Ejercicios07.pdf}{Ejercicios07.pdf}, decidiendo si un conjunto es independiente, si genera el espacio y si puede extraerse una base.
    \item \textbf{Recomendaciones de lectura:} 
    \begin{itemize}
        \item \href{https://alephsub0.org/material-nuevo/acaicedo/determinar-espacio-generado/}{Determinar espacio generado}
        \item \href{https://blog.nekomath.com/algebra-lineal-i-combinaciones-lineales/}{Combinaciones lineales}
        \item \href{https://blog.nekomath.com/algebra-lineal-i-conjuntos-generadores-independencia-lineal-y-bases/}{Conjuntos generadores e independencia lineal}
        \item \href{https://blog.nekomath.com/algebra-lineal-i-problemas-de-combinaciones-lineales-generadores-e-independientes/}{Problemas de combinaciones lineales, generadores e independientes}
        \item \href{https://blog.nekomath.com/algebra-lineal-i-bases-y-dimension-de-espacios-vectoriales/}{Bases y dimensión de espacios vectoriales}
        \item \href{https://blog.nekomath.com/algebra-lineal-i-problemas-de-bases-y-dimension-de-espacios-vectoriales/}{Problemas de bases y dimensión de espacios vectoriales}
    \end{itemize}
    \item \textbf{Visualización de video:} \href{https://youtu.be/TgKwz5Ikpc8?si=7z5e-PTXFkTK7szh}{Espacios vectoriales abstractos}
\end{itemize}

\paragraph{Verificación de aprendizaje:}
\begin{enumerate}[leftmargin=*]
    \item ¿Qué significa que un conjunto de vectores sea linealmente dependiente?
    \item ¿Cuál es la relación entre base, conjunto generador e independencia?
    \item ¿Cómo se determina la dimensión de un espacio vectorial?
\end{enumerate}

%%%%%%%%%%%%%%%%%%%%%%%%%%%%%%%%%%%%%%%%
\section*{Cierre}
%%%%%%%%%%%%%%%%%%%%%%%%%%%%%%%%%%%%%%%%

\paragraph{Tarea:}
Resolver del libro \href{https://catalogobiblioteca.puce.edu.ec/cgi-bin/koha/opac-detail.pl?biblionumber=86081}{Fundamentos de álgebra lineal de Larson}: sección 4.4, ejercicios 2, 5, 7, 11, 15, 21, 22, 29, 30, 31, 32, 54; sección 4.5, ejercicios 7, 8, 9, 10, 11, 12, 13, 21, 23, 25, 27, 33, 35, 41, 43.

\paragraph{Pregunta de investigación:}
\begin{enumerate}[leftmargin=*]
    \item ¿Cómo se puede decidir de forma rápida si un conjunto finito genera un espacio vectorial?
    \item ¿Qué estrategias permiten construir una base a partir de un conjunto generador?
    \item ¿Cómo cambia la noción de independencia lineal cuando se trabaja con conjuntos infinitos?
\end{enumerate}

\paragraph{Para la próxima clase:}  
Visualizar el video \href{https://youtu.be/LYlaRDsi_T8?si=3zNWorDZfwZVS3Eq}{Cambio de Bases} y realizar la lectura de la sección 4.7 del libro de Larson.



\end{document}
